\documentclass[../main.tex]{subfiles}
\begin{document}
\chapter{Python Implementations}
\label{app:python-implementations}

The listings in this appendix give the numerical kernels used by the worked examples. The notebooks contain the complete symbolic derivations, parameter definitions, residual checks, plots, and run instructions; they are the canonical executable artifacts. The snippets below are intentionally limited to the repeated solve steps, so the derivation and validation context is not buried in a long code transcription.

\section{Double-pendulum acceleration}
The notebook derives callable forms of the mass matrix and the vector \(h\) from the symbolic Euler--Lagrange residual. The numerical right-hand side solves the linear system directly rather than forming an inverse.

\begin{lstlisting}[caption={First-order right-hand side for the unforced double pendulum},label={lst:double-pendulum-rhs}]
import numpy as np

def double_pendulum_rhs(time, state, mass_matrix, bias):
	q = state[:2]
	qdot = state[2:]
	matrix = mass_matrix(q)
	h = bias(q, qdot)
	qddot = np.linalg.solve(matrix, -h)
	return np.concatenate((qdot, qddot))
\end{lstlisting}

For applied generalized torques, replace \texttt{-h} by \texttt{torque(time, q, qdot) - h}. The notebook's symbolic check verifies the decomposition before the numerical functions are created.

\section{Five-link constrained acceleration}
For the five-link model, the mass matrix is diagonal in center-of-mass coordinates. The constraint model supplies \(\Phi\), \(J\), and \(\dot J\dot q\). The following kernel adds Baumgarte feedback only to the lower KKT block; set both gains to zero to recover the ideal acceleration constraint at a consistent state.

\begin{lstlisting}[caption={KKT acceleration and reaction-force solve},label={lst:five-link-kkt}]
import numpy as np

def constrained_acceleration(q, qdot, time, model,
							 zeta=1.0, omega=20.0):
	phi, jac = model.constraints_and_jacobian(q)
	bias = model.jdot_qdot(q, qdot)
	stabilization = (
		2.0 * zeta * omega * (jac @ qdot)
		+ omega**2 * phi
	)
	constraint_rhs = -bias - stabilization

	count = jac.shape[0]
	zero = np.zeros((count, count))
	kkt = np.block([
		[model.mass_matrix, jac.T],
		[jac, zero],
	])
	rhs = np.concatenate((
		model.generalized_forces(q, qdot, time),
		constraint_rhs,
	))
	solution = np.linalg.solve(kkt, rhs)
	split = q.size
	return solution[:split], solution[split:]
\end{lstlisting}

The multiplier signs correspond to the chapter convention \(\Phi_{ab}=r_a-r_b\) and \(M\ddot q+J^T\lambda=Q\). A production implementation should also report Jacobian rank and scaled linear-solve residuals and should handle rank loss explicitly, as the canonical notebook does.

\section{State derivative and integration}
The stabilized acceleration kernel defines an ODE for \(y=[q,\dot q]^T\). The code below shows the integrator call used by the five-link example; initial position and velocity projection are performed before this call.

\begin{lstlisting}[caption={State derivative and adaptive integration},label={lst:five-link-integration}]
from scipy.integrate import solve_ivp

def state_derivative(time, state, model, coordinate_count):
	q = state[:coordinate_count]
	qdot = state[coordinate_count:]
	qddot, _ = constrained_acceleration(
		q, qdot, time, model
	)
	return np.concatenate((qdot, qddot))

time_samples = np.linspace(0.0, 10.0, 2001)
initial_state = np.concatenate((q0, qdot0))
solution = solve_ivp(
	lambda time, state: state_derivative(
		time, state, system, q0.size
	),
	(0.0, 10.0),
	initial_state,
	t_eval=time_samples,
	rtol=1e-8,
	atol=1e-10,
	max_step=0.01,
)
if not solution.success:
	raise RuntimeError(solution.message)
\end{lstlisting}

After integration, evaluate \(\|\Phi(q)\|_\infty\), \(\|J(q)\dot q\|_\infty\), and the grounded-pose error at each saved state. Solver success alone does not establish constraint preservation or physical validity.
\end{document}